\section{Number of Islands}
\label{sec:graph-number-of-islands}

\problemstatement{Given a 2D grid of \texttt{'1'} (land) and
\texttt{'0'} (water), count the number of \textbf{islands} -- maximal
groups of land cells connected $4$-directionally.}

\begin{patternbox}
This closing problem of the chapter is
Section~\ref{sec:graph-number-of-provinces}'s connected-component
counting loop, applied to an \textbf{implicit} grid graph instead of an
explicit adjacency matrix -- a land cell's neighbors are simply its
(up to) $4$ adjacent grid cells, with no adjacency list ever
constructed.
\end{patternbox}

\subsection*{Algorithm}

\begin{enumerate}
  \item Maintain a $\textit{visited}$ grid (or mutate the input grid
        directly, marking visited land as \texttt{'0'}) and an
        $\textit{islands}$ counter.
  \item For every cell: if it is land and unvisited, increment
        $\textit{islands}$, then DFS (or BFS) from it, marking every
        $4$-directionally connected land cell visited.
  \item Return $\textit{islands}$.
\end{enumerate}

\subsection*{C++ implementation}

\begin{lstlisting}
#include <bits/stdc++.h>
using namespace std;

void dfs(int r, int c, vector<vector<char>>& grid, vector<vector<bool>>& visited) {
    int rows = grid.size(), cols = grid[0].size();
    if (r < 0 || r >= rows || c < 0 || c >= cols) return;
    if (visited[r][c] || grid[r][c] == '0') return;

    visited[r][c] = true;
    dfs(r - 1, c, grid, visited);
    dfs(r + 1, c, grid, visited);
    dfs(r, c - 1, grid, visited);
    dfs(r, c + 1, grid, visited);
}

int numIslands(vector<vector<char>>& grid) {
    int rows = grid.size(), cols = grid[0].size();
    vector<vector<bool>> visited(rows, vector<bool>(cols, false));
    int islands = 0;

    for (int r = 0; r < rows; r++) {
        for (int c = 0; c < cols; c++) {
            if (grid[r][c] == '1' && !visited[r][c]) {
                islands++;
                dfs(r, c, grid, visited);
            }
        }
    }
    return islands;
}

int main() {
    vector<vector<char>> grid = {
        {'1', '1', '0', '0'},
        {'1', '1', '0', '0'},
        {'0', '0', '1', '0'},
        {'0', '0', '0', '1'}
    };
    cout << numIslands(grid) << endl; // 3
    return 0;
}
\end{lstlisting}

\subsection*{Dry run}

\dryrun{Take the $4\times4$ grid above.}

$(0,0)$: land, unvisited, $\textit{islands}=1$; DFS visits the
top-left $2\times2$ block of land. $(0,1),(1,0),(1,1)$: already
visited, skip. $(2,2)$: land, unvisited, $\textit{islands}=2$; DFS
visits just itself (no land neighbors). $(3,3)$: land, unvisited,
$\textit{islands}=3$; DFS visits just itself. Final answer: $3$
islands.

\begin{complexitybox}
\textbf{Time Complexity.} $O(\textit{rows} \times \textit{cols})$ --
every cell is visited at most once across all DFS calls combined.
\textbf{Space Complexity.} $O(\textit{rows} \times \textit{cols})$ for
the \textit{visited} grid and, in the worst case, the recursion stack
(an entirely land-filled grid).
\end{complexitybox}

\begin{takeawaybox}
This chapter's eight problems form a complete toolkit built from just
two primitives -- DFS and BFS -- applied to three settings (an explicit
adjacency list, an explicit adjacency matrix, and an implicit grid):
plain traversal, traversal with a hash map for building a copy,
component-counting via repeated traversal from unvisited starts, and
multi-source BFS for simultaneous spreading. Every later chapter in
this book's graph coverage layers a new constraint (direction, weight,
ordering) on top of exactly these same two traversal primitives.
\end{takeawaybox}
